\documentclass[11pt]{article}
\usepackage{amsmath,amssymb,geometry}
\geometry{margin=1in}
\pagestyle{plain}
\numberwithin{equation}{section}
\newcommand{\eps}{\varepsilon}
\newcommand{\ZZ}{\mathbb{Z}}

\title{The Gravitational Dyon Pair:\\ Quarter Quantisation on the Poincar\'e Cone}
\author{Jeremy Ebert}
\date{October 2026}

\begin{document}
\maketitle

\begin{abstract}
A mass $m$ and a NUT charge $l$ form gravity's dyon pair: the monopole of the
gravitoelectric field and the monopole of the gravitomagnetic field. That the
pair obeys a Dirac-type quantisation condition has been known in essence
since Misner and Dowker, and is usually left as a remark. This note pins the
condition down. From the $\mathrm{SU}(2)$ Killing triad of Taub--NUT we show
that a test particle's conserved angular momentum satisfies
$\mathbf{J}\cdot\hat{\mathbf{n}}=-2\eps l$ exactly, so the pair's stored
angular momentum is $2\eps l$, twice the electromagnetic value at the same
formalism; the Misner time period $8\pi l$ gives the same lattice by an
independent route. The pair product is therefore quantised in quarters,
$\mu c l/\hbar\in\tfrac14\ZZ$, i.e.\ $\alpha_G(\mu)\,\alpha_G(\tilde\mu)
=(k/4)^2$ with minimum $1/16$, against the electromagnetic
$\alpha\,\alpha_m=(1/2)^2=1/4$. We then exhibit structures of the pair
that appear not to have been stated: its seats in the parabolic separation of
the MIC--Kepler problem are $k/2\pm\eps l$; the monopole enters the spectrum
only through the offset $\Lambda+\tfrac12=\sqrt{(j+\tfrac12)^2-s^2}$ of the
ladder's diameter--secant split; the electromagnetic half-lattice
is the even sublattice of the gravitational quarter-lattice; and a boost
decomposes the state's Casimir between its orbital and pair legs, saturating
at a finite rapidity $\gamma_{\max}=2\sqrt{j(j+1)}/n$. All identities were
verified by exact symbolic computation.
\end{abstract}

\section{Two dyon pairs}

Electromagnetism has one observed monopole, the electric charge, and one
exact zero, $\nabla\!\cdot\!\mathbf{B}=0$. Gravity repeats the pattern one
column over: mass is the observed monopole of the gravitoelectric field, and
the monopole of the gravitomagnetic field---the NUT charge, or dual
mass---is unobserved. In each theory the two charges form a dyon pair, and
in each theory the pair is constrained by single-valuedness.

For the electromagnetic pair the constraint is Dirac's \cite{Dirac}. Its
sharpest derivation is Poincar\'e's: the field of a charge--monopole pair
stores angular momentum $eg/c$, directed along the line joining them,
independently of their separation. The total angular momentum of a test
charge therefore has a fixed radial component,
\begin{equation}
\mathbf{J}\cdot\hat{\mathbf{n}}=-eg,
\qquad
J^2=L^2+(eg)^2,
\end{equation}
with $L\perp\hat{\mathbf{n}}$ the orbital part; the position direction
$\hat{\mathbf{n}}$ is confined to a cone about $\mathbf{J}$ (the Poincar\'e
cone), and $\mathrm{SU}(2)$ consistency forces $eg\in\tfrac{\hbar}{2}\ZZ$.
In dimensionless couplings, with $\alpha=e^2/\hbar c$ and
$\alpha_m=g^2/\hbar c$,
\begin{equation}
\alpha\,\alpha_m=\Big(\frac{eg}{\hbar c}\Big)^2=\Big(\frac{k}{2}\Big)^2,
\qquad k\in\ZZ,
\end{equation}
with minimum $1/4$: the pair product is the square of the pair's
angular-momentum quantum, $\hbar/2$.

The gravitational analogue---Taub--NUT \cite{Taub,NTU} read as a
gravitational dyon---goes back to Demia\'nski--Newman and to Dowker
\cite{Dowker}, and the existence of a quantisation condition relating mass
and NUT charge, obtained from the periodicity of the time coordinate, is
recorded e.g.\ in the double-copy literature \cite{LMOW} following Misner
\cite{Misner}. What seems not to be in print is the condition's coefficient
in pair-product form, nor the structures of Sections 5--8. We supply both
routes to the coefficient here because their agreement is the content: the
factor of two that decides between $1/4$ and $1/16$ is the fibration's own
$2$, read once from particle dynamics and once from topology.

\section{The pair term from the Killing triad}

Write Taub--NUT in the standard parametrisation, in geometric units, with
$\Sigma=r^2+l^2$ and $f(r)=(r^2-2mr-l^2)/\Sigma$:
\begin{equation}
ds^2=-f\,\chi^2+\frac{dr^2}{f}+\Sigma\,(d\theta^2+\sin^2\theta\,d\varphi^2),
\qquad
\chi=dt+2l\cos\theta\,d\varphi .
\end{equation}
The isometry group is $\mathrm{SU}(2)\times\mathrm{U}(1)$. The
$\mathrm{SU}(2)$ factor is not the naive rotation group of the
$(\theta,\varphi)$ sphere: each would-be rotation must be accompanied by a
$t$-shift to remain a symmetry. Solving the Killing equation with $f(r)$
left arbitrary, the shift is forced (the opposite sign fails), giving the
triad
\begin{align}
\xi_x&=2l\,\frac{\cos\varphi}{\sin\theta}\,\partial_t
-\sin\varphi\,\partial_\theta-\cos\varphi\cot\theta\,\partial_\varphi,\\
\xi_y&=2l\,\frac{\sin\varphi}{\sin\theta}\,\partial_t
+\cos\varphi\,\partial_\theta-\sin\varphi\cot\theta\,\partial_\varphi,\\
\xi_z&=\partial_\varphi ,
\end{align}
which closes as an $\mathfrak{su}(2)$ algebra,
$[\xi_i,\xi_j]=-\epsilon_{ijk}\xi_k$. The closure is a real constraint: the
alternative third generator $\partial_\varphi+2l\,\partial_t$ is also a
Killing vector, being a linear combination of symmetries, but it does
\emph{not} close the algebra with $\xi_x,\xi_y$.

For a unit-mass geodesic let $\eps=-p_t$ (specific energy) and
$k=p_\varphi$. The conserved charges $Q_i=\xi_i\cdot p$ assemble into a
total angular momentum $\mathbf{J}$ satisfying two exact identities:
\begin{equation}
\mathbf{J}\cdot\hat{\mathbf{n}}=-2\eps l,
\qquad
J^2=\Sigma^2\dot\theta^{\,2}
+\frac{\big(k+2\eps l\cos\theta\big)^2}{\sin^2\theta}
+\big(2\eps l\big)^2 .
\end{equation}
This is Poincar\'e's structure (1.1) with the replacement
\begin{equation}
eg\ \longrightarrow\ 2\eps l .
\end{equation}
Three remarks. First, the ``electric'' member of the gravitational pair is
the particle's \emph{energy}: the gravitomagnetic charge couples to $\eps$,
as the Kaluza--Klein reading of the $t$-fibration requires. At rest at
infinity $\eps=1$. Second, $\mathbf{J}\cdot\hat{\mathbf{n}}$ fixed confines
$\hat{\mathbf{n}}$ to a cone about $\mathbf{J}$ with
$\cos\alpha=-2\eps l/|\mathbf{J}|$: the Poincar\'e cone, in the
gravitational theory, with the pair term as one leg of the cone relation
$J^2=(\text{orbital})^2+(\text{pair})^2$. Third, the coefficient is $2$,
not $1$: it is the $2$ in $\chi=dt+2l\cos\theta\,d\varphi$, and it is
physical, since the same $l$ appears in $f$ and $\Sigma$ and admits no
reparametrisation.

\section{The Misner route}

The same coefficient controls the string. Near the poles,
$\chi\approx dt\pm2l\,d\varphi$; the regular time coordinates
$t_N=t+2l\varphi$ and $t_S=t-2l\varphi$ disagree by $4l\varphi$ on the
overlap, hence by $8\pi l$ after one wind. Removing the Misner string
therefore requires $t$ periodic with period $T_t=8\pi l/c$ \cite{Misner},
and single-valuedness of $e^{-iEt/\hbar}$ quantises energy in steps
\begin{equation}
E_k=\frac{2\pi\hbar}{T_t}\,k=\frac{\hbar c}{4l}\,k,\qquad k\in\ZZ .
\end{equation}
Setting $E=\mu c^2$ places a test mass on the ladder. Both routes give the
pair lattice of the next section, and both read the answer off the same
$2$: Route~A as the strength of the fibration's monopole potential in the
particle's angular momentum, Route~B as the Misner period $8\pi l=4\pi\cdot2l$.

\section{The pair product: quarters}

Since $\mathbf{J}\cdot\hat{\mathbf{n}}$ is a component of an
$\mathrm{SU}(2)$ generator and the remaining part of $\mathbf{J}$ is
orthogonal to $\hat{\mathbf{n}}$, equation (2.5) with $\eps=1$ requires
$2\mu_{\mathrm{geo}} l_{\mathrm{geo}}/\ell_P^2\in\tfrac12\ZZ$ for a test
mass $\mu$, i.e.
\begin{equation}
\frac{\mu_{\mathrm{geo}}\,l_{\mathrm{geo}}}{\ell_P^2}\in\frac14\ZZ
\qquad\Longleftrightarrow\qquad
\frac{\mu\, c\, l}{\hbar}\in\frac14\ZZ ,
\end{equation}
the second form showing that $G$ cancels identically: the condition is
dimensionless, $G$-free, and a property of the pair alone. With the dual
mass $\tilde\mu=c^2l/G$ and the dimensionless gravitational couplings
$\alpha_G(\mu)=G\mu^2/\hbar c$, $\alpha_G(\tilde\mu)=G\tilde\mu^2/\hbar c$,
\begin{equation}
\boxed{\ \alpha_G(\mu)\,\alpha_G(\tilde\mu)=\Big(\frac{k}{4}\Big)^2,\qquad
k\in\ZZ,\qquad \text{minimum } \frac{1}{16}.\ }
\end{equation}
Gravity's pair quantum is $\hbar/4$, against electromagnetism's $\hbar/2$.

The exponent structure is worth displaying. The electromagnetic product is
$(1/2)^2$: one factor $1/2$ from single-valuedness on the half-ladder,
squared by the pairing. The gravitational product is
$1/16=(1/2)^{2\cdot2}$: the same half, halved once more by the spin-2
fibration coefficient, then squared by the pairing. The quarter is also the
cone's native quarter: on the angular-momentum ladder the Casimir obeys
$T^2=j(j+1)+\tfrac14$ with $T=j+\tfrac12$, i.e.\ the offset between the
arithmetic radius $j+\tfrac12$ and the geometric Casimir $\sqrt{j(j+1)}$ is
$1/2$, and its square---the deficit in $T^2$---is $1/4$. The pair quantum
$\hbar/4$ and the Casimir deficit $1/4$ are the same quarter seen in two
places: one in a product of charges, one in a spectrum. The mechanism has
been isolated before \cite{EbertCas} as \emph{dyadic quadratic
normalisation}: a canonical half-coordinate becomes a quarter under a
quadratic operation, with four exact realisations (the factor-cell
midpoint square, the shifted compact and noncompact Casimirs, the
primitive null-edge midpoint norm, and the cyclotomic relative-norm
normalisation). The pair product (4.2) is that mechanism at dyadic depth
two---half to quarter for the pair quantum, quarter to sixteenth for its
square---and its first realisation in a product of charges. The same
square root governs the power laws on the parabolas themselves: a
parabolic segment has area $\propto h^{3/2}$ with $3/2=1+\tfrac12$, and the
bouncer spectrum warps as $E\propto(n-\tfrac14)^{2/3}$ with
$2/3=1/(1+\tfrac12)$ \cite{EbertSecant}.

\section{Parabolic seats}

The Kepler problem separates in parabolic coordinates
$\xi=r(1+\cos\theta)$, $\eta=r(1-\cos\theta)$, and its parabolic quantum
numbers are the seats from which the hydrogenic fine-structure bookkeeping
is read off. The monopole deformation of Kepler---the MIC problem
\cite{MIC,Zwanziger}---is the flat-space twin of Taub--NUT geodesic motion.
Its Hamilton--Jacobi equation separates in the same coordinates. Exactly
(including the Zwanziger term $+s^2/r^2$, without which a remainder
$-s^2/(\xi+\eta)$ obstructs separation), with $s$ the Poincar\'e strength
and $\kappa$ the Coulomb coupling:
\begin{equation}
\Big[\xi U'^2+\frac{((k-s)/2)^2}{\xi}-\frac{E\xi}{2}\Big]
+\Big[\eta V'^2+\frac{((k+s)/2)^2}{\eta}-\frac{E\eta}{2}\Big]=\kappa .
\end{equation}
Each parabola carries one charge, $(k\mp s)/2$: centre $k/2$, offset $s/2$.
Hydrogen is $s=0$, both seats equal to $(k/2)^2$.

Taub--NUT's Hamilton--Jacobi equation has MIC's angular and pair sectors
\emph{identically}, at $s=2\eps l$: subtracting the two equations,
\begin{equation}
\mathrm{TN}-\mathrm{MIC}=-(2mr+l^2)\,p_r^2-(2\eps l)^2 ,
\end{equation}
which is purely radial---a $p_r^2$ weight and a constant, the sector
absorbed by the standard time reparametrisation identifying the two flows
\cite{IK}. The parabolic content is undeformed. Hence Taub--NUT's
parabolic seats are
\begin{equation}
\frac{k}{2}-\eps l,\qquad \frac{k}{2}+\eps l :
\end{equation}
the diameter is $k$, the arithmetic centre is $k/2$, and the offset is the
pair product $\eps l$ itself. The twist does not build a new lattice; it
splits each hydrogen seat symmetrically about its centre by the pair term
of Section~2. The state lattice of the separation thus runs in halves,
while the pair lattice of Section~4 runs in quarters---one additional
power of two of resolution, supplied by the fibration.

\section{The monopole is the offset}

The identity (2.5) says $J^2=L^2+s^2$ with $s$ the pair strength ($s=eg$
electromagnetically, $s=2\eps l$ gravitationally): the orbital part is what
remains of the Casimir after the pair takes its leg. Quantising, the
mechanical angular momentum of the MIC system has eigenvalues
$\Lambda(\Lambda+1)$ with \cite{MIC}
\begin{equation}
\Lambda(\Lambda+1)=j(j+1)-s^2,
\qquad\text{i.e.}\qquad
\Lambda+\frac12=\sqrt{\Big(j+\frac12\Big)^2-s^2}\;,
\end{equation}
and the bound spectrum is hydrogenic in the shifted rung,
$E=-\kappa^2/(2N^2)$ with $N=n_r+\Lambda+1$, recovering hydrogen at $s=0$.

Equation (6.1) is a diameter--secant split \cite{EbertSecant}: the
arithmetic mean is $j+\tfrac12$, the offset is $s$, and the rung the
spectrum actually climbs is the geometric mean
$\sqrt{\mathrm{AM}^2-\mathrm{offset}^2}$. A monopole therefore does not
deform the ladder; it offsets it, entering through the offset slot alone.
The parabolic seats (5.3) show the same entry in the other coordinate
system: centre $k/2$ untouched, offset $\eps l$.

Consequently both monopoles obey the machine's power laws, not by
generating them but by inheriting them: the laws belong to the parabola
and its square root---segment area $\propto h^{3/2}$ with
$3/2=1+\tfrac12$, spectral warp $(n-\tfrac14)^{2/3}$ with
$2/3=1/(1+\tfrac12)$, and the dyadic half-to-quarter normalisation of
Section~4---and the offset is the only door a monopole has into that
machine. In the diameter--secant picture the offset $(n-k)/2$ was always
the free parameter, the position chosen by the secant. The pair gives
that slot a physical occupant: the product of the pair's two charges. The
question ``what fixes the offset'' and the question ``what carries
monopole charge, and how much'' are the same question.

\section{Intertwining of the two lattices}

Equations (1.2) and (4.2) place the two theories' pair lattices on one
grid:
\begin{equation}
\frac{eg}{\hbar c}\in\frac12\ZZ\subset\frac14\ZZ\ni\frac{\mu c l}{\hbar}.
\end{equation}
The electromagnetic half-lattice is the \emph{even sublattice} of the
gravitational quarter-lattice: the minimal electromagnetic pair quantum
$\hbar/2$ is the $k=2$ rung of the gravitational lattice, $2\times\hbar/4$.
Electromagnetism occupies every second rung; the odd rungs belong to the
twist sector alone. Two lattices, one grid---and the grid is fixed by
single-valuedness in both theories, not fitted in either.

\section{Boost decomposition}

A boost along the symmetry axis acts on the pair through $\eps$ alone:
$k$, $l$ and $|\mathbf{J}|$ are unchanged, while a boost from rest sends
$\eps\to\cosh\psi$. The parabolic seats (5.3) split hyperbolically about
their fixed centre,
\begin{equation}
\text{seats}=\frac{k}{2}\pm l\cosh\psi ,
\end{equation}
and the Casimir budget $J^2=(\text{orbital})^2+(2\eps l)^2$ is
transferred, at fixed $J^2$, from the orbital leg to the pair leg: the
boost decomposes the state into the pair. Writing the pair fraction with
both quantisations imposed, $\mu c l/\hbar=n/4$ and
$|\mathbf{J}|^2/\hbar^2=j(j+1)$, the transfer saturates when the orbital
leg vanishes, $2\eps l=|\mathbf{J}|/\mu$, at the finite rapidity
\begin{equation}
\gamma_{\max}=\frac{2\sqrt{j(j+1)}}{n},
\qquad
\psi_{\max}=\operatorname{arcosh}\gamma_{\max}.
\end{equation}
At the saturation the Poincar\'e cone has closed,
$\hat{\mathbf{n}}$ locked antiparallel to $\mathbf{J}$. The minimal state
$(j,n)=(\tfrac12,1)$ has $\gamma_{\max}=\sqrt3$,
$\psi_{\max}=\operatorname{arcosh}\sqrt3\approx1.146$, and rest-frame split
pair $1/3$ / orbital $2/3$ (arithmetic of the minimal rung, recorded as
such). Finiteness cuts the other way as well: a shell $j$ admits at rest
only the pair charges $n\le2\sqrt{j(j+1)}$---$n=1$ for $j=\tfrac12$;
$n=1,2$ for $j=1$; $n=1,2,3$ for $j=\tfrac32$.

On the source side, the $\psi\to\infty$ image of a boost is the
Aichelburg--Sexl shockwave \cite{AS}: a fully boosted field collapses onto
the null cone as an impulsive wave. The pair-side saturation (8.2) is the
finite, quantised counterpart for a test state: the cone closes before
the boost diverges.

\section{The action is the integration}

None of the preceding structure floats free of an action: the action is
the Hamilton--Jacobi integral itself. The characteristic function
$W=U(\xi)+V(\eta)$ of Section~5 \emph{is} Hamilton's action for the pair
system, separated into its two parabolic pieces, and each quantisation
encountered above is a statement about its periods. Route~A fixes the
angular period: $\mathbf{J}$ is the action variable of the
$\mathrm{SU}(2)$ motion, and the pair term is the part of that period
carried by the twist. Route~B fixes the time period: single-valuedness of
$e^{-iEt/\hbar}$ on the Misner cycle. The parabolic power laws are
periods again---the Archimedean segment area $\propto h^{3/2}$ is the
action of the linear-potential (bouncer) problem, whose Bohr--Sommerfeld
periods, with their Maslov quarter, give the warp
$E\propto(n-\tfrac14)^{2/3}$. The lattice of Section~4 is, in one phrase,
the period lattice of the pair's action in units of $\hbar$: halves where
the action is electromagnetic, quarters where it is gravitational.

This places the structure one rung below a field theory, not outside the
theory business altogether. A worldline effective theory of the pair is
built on precisely this action, and for such a theory the results above
are exact period data---matching conditions it must reproduce: the
quarter lattice (4.2), the seats (5.3) and the offset law (6.1), and the
boost saturation (8.2). What is absent is the field layer: a Lagrangian
density for the gravitoelectric and gravitomagnetic fields jointly, with
an expansion organised on top of the worldline action. Section~10 returns
to what that absence does and does not affect.

\section{Scope}

Nothing here detects the pair. NUT charge is unobserved; Taub--NUT without
the periodic-time identification carries the Misner string, and with it,
closed timelike curves. Equations (4.2), (5.3) and (8.2) are statements
about the solution space of general relativity and the representation
theory of its isometries, not about nature's inventory. Nor does the pair
product derive $G$: $l$ is a free charge, and the lattice constrains
products, not values---exactly Dirac's situation, where the condition
quantises $eg$ given a monopole and leaves $\alpha$ free.

What the structure does supply is the \emph{form} a gravitational
selection principle would take if one exists in this family: dimensionless,
$G$-free once written as (4.1), and pair-dependent at once, with its
symmetry---the boost action of Section~8---already realised in the
gravitomagnetic sector that frame-dragging measurements probe. The missing
ingredient is the same one Dirac's condition waited on: a physical pair
with one member known. Until then the correct classification is a
quantisation of permitted pairs, on a quarter lattice, with electromagnetism
on its even rungs.

\begin{thebibliography}{99}
\bibitem{Dirac} P.~A.~M. Dirac, Quantised singularities in the
electromagnetic field, \emph{Proc.\ Roy.\ Soc.\ A} \textbf{133} (1931)
60--72.
\bibitem{Taub} A.~H. Taub, Empty space-times admitting a three parameter
group of motions, \emph{Ann.\ of Math.} \textbf{53} (1951) 472--490.
\bibitem{NTU} E.~Newman, L.~Tamburino and T.~Unti, Empty-space
generalization of the Schwarzschild metric, \emph{J.\ Math.\ Phys.}
\textbf{4} (1963) 915--923.
\bibitem{Misner} C.~W. Misner, The flatter regions of Newman, Unti, and
Tamburino's generalized Schwarzschild space, \emph{J.\ Math.\ Phys.}
\textbf{4} (1963) 924--937.
\bibitem{DN} M.~Demia\'nski and E.~T. Newman, A combined Kerr--NUT
solution of the Einstein--Maxwell equations, \emph{Bull.\ Acad.\ Polon.\
Sci.} \textbf{14} (1966) 653--657.
\bibitem{Dowker} J.~S. Dowker, The NUT solution as a gravitational dyon,
\emph{Gen.\ Rel.\ Grav.} \textbf{5} (1974) 603--613.
\bibitem{LMOW} A.~Luna, R.~Monteiro, D.~O'Connell and C.~D. White, The
classical double copy for Taub--NUT spacetime, \emph{JHEP} \textbf{11}
(2015) 143.
\bibitem{MIC} H.~V. McIntosh and A.~Cisneros, Degeneracy in the presence
of a magnetic monopole, \emph{J.\ Math.\ Phys.} \textbf{11} (1970)
896--916.
\bibitem{Zwanziger} D.~Zwanziger, Quantum field theory of particles with
both electric and magnetic charges, \emph{Phys.\ Rev.} \textbf{176}
(1968) 1480--1488.
\bibitem{IK} T.~Iwai and Y.~Katayama, On extended Taub-NUT metrics,
\emph{J.\ Geom.\ Phys.} (1993--94), on the Taub--NUT $\leftrightarrow$
MIC--Kepler correspondence under time reparametrisation.
\bibitem{AS} P.~C. Aichelburg and R.~U. Sexl, On the gravitational field
of a massless particle, \emph{Gen.\ Rel.\ Grav.} \textbf{2} (1971)
303--312.
\bibitem{BH} C.~Bunster, S.~Cnockaert, M.~Henneaux and R.~Portugues,
Monopoles for gravitation and for higher spin fields, \emph{Phys.\ Rev.\
D} \textbf{73} (2006) 105014.
\bibitem{EbertCas} J.~Ebert, Casimir completion and primitive null
diamonds: edge centering, cyclotomic normalization, and the ramified
Boolean return, manuscript (2026).
\bibitem{EbertSecant} J.~Ebert, The diameter--secant machine and the
gravity road, with addendum: the Airy comb and the two powers of the
machine, research notes (2026).
\end{thebibliography}

\end{document}
